\section*{September 21, 2026: Neighborhood sales mix explains the assessment-median spike}
\textit{Drafted by Codex; investigation requested by Jacob.}

More sales near expensive sites explain the spike in treated sales' median
2006 assessment. Holding neighborhood shares fixed nearly removes it, while
the adjusted log-price gap between 2010 and 2012 remains. Independent data
show higher average foreclosure-filing rates at control locations.

Among transactions with usable initial assessments, the treated median
rises from \$16,072 in 2011 to \$26,345 in 2012. These are fixed historical
values of the homes sold, not changes in assessments. Sales near Trumbull,
King, Peabody, Duprey/Humboldt and Near North increase from 28 of 99 to
66 of 134. Exactly half the 2012 sales have assessments above \$25,000;
the middle two observations are \$24,040 and \$28,650. The median is their
average. Five parcels sell twice, with distinct deeds, dates and prices.
Counting each parcel once lowers the 2012 median to \$20,880: repeated
transactions amplify the underlying shift in which homes sell.

For the 14 treated sites with sales in every year from 2008 through 2013,
the observed median rises from \$16,417 to \$29,043. On those exact same
transactions, holding each site's share at its pooled 2008--2011 share changes
the medians to \$16,121 and \$16,371. A symmetric decomposition attributes
\$6,010 (75.5\%) of the \$7,959 increase in the mean assessment to changing
site shares, with the rest coming from different homes sold within sites.
The working transaction weights remain unchanged; this is a composition check.

DePaul IHS publishes community-area foreclosure filings per 100 residential
parcels, constructed from court and assessor records. We preserve its
September 21, 2026 table and Chicago's community boundaries, assign each
sale point to exactly one area, and average rates using fixed 2008--2011
sale-location shares. The resulting rates are 3.85 treated versus 4.76 control
in 2008, and 3.44 versus 4.01 in 2010. Both groups exceed the citywide
2010 rate of 3.0. These are broad area exposures covering all residential
property types, not filing rates inside our circles or on the sold homes.
Filings are not completed foreclosures.

The sales indicators tell a similar, qualified story. Broad bank/lender-related
seller names account for 31.0\% of treated versus 40.3\% of control sales
in 2010; narrow judicial/GSE/HUD names account for 4.3\% versus 14.2\%.
Neither is a verified distress classification. Treated sales below \$50,000
in 2022 dollars fall from 39.3\% to 24.8\% between 2011 and 2012, while
control shares rise from 37.1\% to 40.4\%. Treated building size and townhouse
shares also increase. An explicit corporate-suffix buyer proxy shows no
persistent treated excess over 2010--2012; it does not classify investors.

Changing the log model's reference year barely changes the 2010 point:
$-19.9\%$ relative to 2012 versus $-18.8\%$ relative to 2008. With 2008 as the
reference, the relative decline to 2010 has a 95\% interval of
$[-34.9\%,1.3\%]$. The 2010--2012 relative rebound is 24.9\%, with an
interval of $[8.2\%,44.0\%]$. Rebasing leaves fitted values and the joint
pre-test unchanged. Actually removing 2012 leaves a similar 18.7\% decline;
the remaining pre-test has p=0.187, which does not establish parallel trends.
Seller-category-by-year controls and removal of bank-related sellers leave
substantial movement. On the identical 4,919-sale common-site sample for
2008--2013, the adjusted 2010--2012 rebound remains about 26\% with either
observed or fixed site shares. Omitting each of the 68 sites individually
leaves the 2010-versus-2012 point estimate negative. Slower recovery among
more distressed controls remains plausible, but unproven.

The follow-up on sales volume finds a quiet 2011 followed by a broader
increase: the five selected areas have 39 assessment-matched sales in 2010,
28 in 2011, 66 in 2012 and 65 in 2013. Trumbull alone goes from 21 to 11
to 21 to 15. Removing the assessment requirement leaves a 32-to-68 increase
between 2011 and 2012; before price and county-quality screens the counts
are 34 and 78. Those earlier records still satisfy the geography, property
class, single-parcel and single-card restrictions. These are sample counts,
not turnover rates for the full housing stock.

The 66 sales in 2012 involve 63 parcels and 65 distinct recorded buyer names.
No normalized buyer or seller name appears more than twice. Sales increase
in every quarter relative to 2011. Bank-related seller matches rise from
four to seven, while other known sellers rise from 24 to 59. King accounts
for five of the seven 2012 bank-related matches. Townhouse sales increase
from four to sixteen; near Duprey/Humboldt, the median building age changes
from 121 to 13 years and median size from 968 to 2,082 square feet. These
are changes in which homes sell. Only one of the 66 homes has a recorded
construction year after 2006, and that year is 2007. The names and building
dates do not establish beneficial ownership, investor intent, distress or
renovation histories.

The Illinois Association of REALTORS' January 22, 2013 release reports a
22.4\% increase in Chicago home sales in 2012. Its MLS series includes
condominiums. A broader recovery is consistent with our increase but does
not explain why these particular areas rose so much: assessment-matched
sales at other treated sites decline from 71 to 68. The five sites were
selected after observing their contribution to the assessment shift.
Neither the selection nor the sales records identify owners' reasons for
selling. Source: \href{https://www.prnewswire.com/news-releases/illinois-sees-home-sales-increases-in-december-2012-notches-229-percent-sales-gain-over-2011-187867161.html}{Illinois Association of REALTORS, 2012 year-end release}.

{\small Reproduction: Make in
\texttt{tasks/audits/home\_recession\_composition/code}, working tree based
on \texttt{e76f4dd}. The seven-page packet and saved reports contain source
rates, joins, site contributions, all specification checks and omissions.
Independent OLS and weighted-median calculations reproduce the results.
The separate \texttt{site\_sales} figure and dataset preserve annual sample
counts, quarterly sales, name concentration and property characteristics.
Source: \href{https://www.housingstudies.org/data-portal/browse/?indicator=foreclosures-100-residential-parcel}{DePaul IHS foreclosure series}.}

\clearpage
\begin{center}
\includegraphics[width=\textwidth]{../input/recession_median_spike.png}

{\small Holding neighborhood shares fixed nearly removes the treated median-assessment spike.}
\end{center}
